\documentclass[11pt,a4paper]{article}
\usepackage[utf8]{inputenc}
\usepackage[T1]{fontenc}
\usepackage{lmodern}
\usepackage[margin=2.6cm]{geometry}
\usepackage{amsmath,amssymb}
\usepackage{booktabs}
\usepackage{array}
\usepackage{enumitem}
\usepackage[colorlinks=true,linkcolor=blue,urlcolor=blue]{hyperref}
\setlength{\parindent}{0pt}
\setlength{\parskip}{6pt}

\newcommand{\HS}{\mathrm{HS}}
\newcommand{\LS}{\mathrm{LS}}
\newcommand{\Dq}{Dq}
\newcommand{\kb}{k_{\mathrm{B}}}

\title{\bfseries Memoria del proyecto original\\[2mm]\large Transición de espín (SCO) en Fe(II)}
\author{Alejandro Rodríguez}
\date{Noviembre 2026}

\begin{document}
\maketitle

\begin{quote}
\emph{Documento de reconstrucción personal.} Objetivo: refrescar la memoria de
\emph{qué} se quería hacer, \emph{cómo} se intentó (los códigos originales) y
\emph{qué falló}, para retomar el doctorado con una base sólida. Todo lo que se
afirma de cada código está verificado contra el fuente en \texttt{~/Dropbox/src/}.
\end{quote}

\tableofcontents

\section{Resumen ejecutivo}

\begin{itemize}
\item El proyecto modela la \textbf{transición de espín (SCO)} de Fe(II) en
  sólidos tipo $\mathrm{Fe(py)_2[Ni(CN)_4]}$ (clatrata de Hofmann / análogos de
  azul de Prusia), donde el Fe pasa de \textbf{alto espín} ($\HS$, $S=2$) a
  \textbf{bajo espín} ($\LS$, $S=0$) al bajar la temperatura, con
  \textbf{histéresis} (efecto cooperativo).
\item El mecanismo microscópico es de \textbf{campo de ligandos}: el
  desdoblamiento octaédrico $\Delta_o = 10\,\Dq$ decide si los electrones ocupan
  los orbitales antienlazantes $e_g$ ($\HS$, enlace largo) o solo los $t_{2g}$
  ($\LS$, enlace corto).
\item Se escribieron \textbf{cinco códigos} (2016--2019) en dos filosofías:
  (1) \emph{modelo de red elástica/Ising} (\texttt{project\_sco}, \texttt{one\_D})
  y (2) \emph{modelo de un centro con $\Dq$} (\texttt{prueba}).
\item El código que se recordaba ("un centro, $\Dq$ y los parámetros $A,B$") es
  \texttt{src/prueba/main.cpp}, basado en el \emph{review} de Levchenko et al.\ (2014).
\item Problema transversal a todos: \textbf{unidades} ($\Dq$ vs $\Delta_o$, $\kb$
  en cm$^{-1}$) y varios bugs menores.
\end{itemize}

\section{Propósito de la modelación}

\begin{center}
\fbox{\parbox{0.95\linewidth}{\small
\textbf{El espíritu del proyecto (recogido en \texttt{prueba/main.cpp}):}\\[1mm]
\textit{variar la distancia Fe--N y ver cómo es el \textbf{campo de ligandos}
$\Dq(r)$ el que provoca, en su raíz, la transición de espín
$\HS\leftrightarrow\LS$. Sobre esta base se añadirán luego los factores que
inciden sobre el campo de ligandos (temperatura, presión, cooperatividad de
red, sustitución de ligandos).}
}}
\end{center}

Es decir: el modelo parte de \emph{un solo centro} y de la relación más
fundamental del problema --- la dependencia del campo de ligandos con la
geometría, $\Dq(r)$ --- y de ahí construye todo lo demás. Esta es la misma idea
que implementaba \texttt{prueba/main.cpp} (energías $\HS$/$\LS$ lineales en
$\Dq$, con el cruce de Tanabe--Sugano) y que ahora recupera el paquete
\texttt{sco0} de forma corregida (ver \texttt{GUIA\_PARA\_CONTINUAR.pdf}).

\section{Contexto y objetivo original}

\textbf{Doctorado (2016--2019).} La parte experimental consistió en sintetizar y
medir complejos de Fe(II) con piridinas sustituidas ($3{,}5$-FBrPy,
$3{,}5$-ClFPy, $2$-FPy, \dots) y caracterizarlos magnéticamente (ZFCFC, $M$ vs $H$,
susceptibilidad AC) y espectroscópicamente (Raman/IR). La parte teórica (los
códigos) buscaba un modelo mínimo que reprodujera la \textbf{curva de conversión
$\HS\to\LS$ en función de la temperatura} y su \textbf{histéresis}.

En la tesis (Cap.~3) se describe así: la transición $\LS\leftrightarrow\HS$ es la
promoción de electrones de los orbitales $e_g$ (antienlazantes) a los $t_{2g}$
(no enlazantes); el estímulo térmico induce un cambio del gap $\Dq$ ($\Delta_o$),
de modo que algunos niveles se vuelven accesibles y otros no. Es decir: la
temperatura \emph{cambia el campo de ligandos efectivo} $\Dq$ (vía la
dilatación/contracción de la red), y eso dispara el cruce $\HS\leftrightarrow\LS$.

\section{Fundamentos físicos}

\subsection{El ion $d^6$ en campo octaédrico}

En un entorno octaédrico los 5 orbitales $d$ se desdoblan en dos grupos:

\begin{center}
\begin{tabular}{lcl}
$e_g$ (antienlazante) & $+6\,\Dq$ & (2 orbitales) \\
$t_{2g}$ (no enlazante) & $-4\,\Dq$ & (3 orbitales)
\end{tabular}
\end{center}

\begin{itemize}
\item \textbf{$\HS$} (alto espín, $S=2$, término ${}^5T_{2g}$):
  configuración $(t_{2g})^4(e_g)^2$. Dos electrones en $e_g$ $\Rightarrow$
  enlace Fe--N \emph{largo}.
\item \textbf{$\LS$} (bajo espín, $S=0$, término ${}^1A_{1g}$):
  configuración $(t_{2g})^6$. Sin electrones en $e_g$ $\Rightarrow$ enlace
  Fe--N \emph{corto} (contracción de red).
\end{itemize}

Energía orbital (sin repulsión electrónica):
\begin{align}
E(\HS) &= 4\,(-4\Dq) + 2\,(+6\Dq) = -4\,\Dq\\
E(\LS) &= 6\,(-4\Dq) = -24\,\Dq
\end{align}

Por eso $\LS$ gana al crecer $\Dq$. Para que $\HS$ sea estable a $\Dq$ pequeño hay
que añadir la \textbf{energía de apareamiento} $P$ (poner dos electrones en el
mismo orbital $t_{2g}$ cuesta energía). El cruce ocurre cuando:
\begin{equation}
-4\Dq = -24\Dq + P \;\Rightarrow\; 20\,\Dq = P \;\Rightarrow\; \Dq_{\mathrm{cruce}} = P/20 .
\end{equation}

\subsection{Diagrama de Tanabe--Sugano y parámetros de Racah}

Los niveles reales (con repulsión $e^-$--$e^-$) se expresan en función de
$\Dq/B$ con los \textbf{parámetros de Racah} $A$, $B$, $C$:
\begin{itemize}
\item El término ${}^5T_{2g}$ ($\HS$) es \textbf{independiente} de $B$ y $C$
  (solo depende de $\Dq$), por eso la recta $-4\Dq$ es exacta.
\item El término ${}^1A_{1g}$ ($\LS$) sube con $B$ y $C$; su cruce con
  ${}^5T_{2g}$ ocurre cerca de $\Delta_o \approx 20B$ (para $C/B \approx 4.4$).
\item Para Fe(II) libre, $B\approx 1058$ cm$^{-1}$ ($\approx 917$ cm$^{-1}$ en
  complejos) y $C\approx 4B$. El parámetro $A$ es un corrimiento promedio (no
  afecta las diferencias $\HS$--$\LS$).
\end{itemize}

En \texttt{prueba/main.cpp} se \textbf{linealizó} esto: $E_{\HS}=-4\Dq$,
$E_{\LS}=-24\Dq+P$, y se hizo Metropolis en el espacio de $\Dq$.

\subsection{Cooperatividad e histéresis}

En el sólido los centros no son independientes: un centro $\HS$ es más grande y
deforma la red, favoreciendo que los vecinos también pasen a $\HS$. Esa
realimentación elástica hace la transición \textbf{cooperativa} y de
\textbf{primer orden} $\Rightarrow$ aparece un \textbf{lazo de histéresis} al
calentar y enfriar. \texttt{project\_sco} y \texttt{one\_D} implementan esa
cooperatividad con resortes entre vecinos; \texttt{prueba} (un centro) no puede
tenerla.

\section{Cronología de los códigos}

\begin{center}
\begin{tabular}{lllll}
\toprule
Código & Fecha & Lenguaje & Qué hace & Estado \\
\midrule
\texttt{DMF.beta.*} & 2016--18 & C++ & Procesado de magnetometría (MvsH, ZFCFC, AC) & Funcional \\
\texttt{one\_D} & Nov 2018 & C++ & Cadena SCO 1D (40 centros), ciclo térmico & Funciona \\
\texttt{project\_sco} & Dic 2018 & C++/CMake & Red 3D elástica/Ising, Metropolis + NPT & Incompleto \\
\textbf{\texttt{prueba}} & \textbf{Mar 2019} & C++ & \textbf{Un centro, $\Dq$ + $B$ (Racah), Metropolis} & Bugs de unidades \\
\texttt{project\_one} & 2019--20 & C++/Qt & GUI de visualización & Prototipo \\
\bottomrule
\end{tabular}
\end{center}

\section{\texttt{prueba/main.cpp}: el código de un centro ($\Dq$ + $B$)}

Este es el código buscado. Ruta: \texttt{~/Dropbox/src/prueba/main.cpp}
(git: ``first version'', 18-mar-2019).

\subsection{Qué hace}

Un único átomo \texttt{hierro} con espín booleano \texttt{spin\_b}
($1=\HS$, $0=\LS$), y una energía $\HS$/$\LS$ \textbf{lineal en $\Dq$}:

\begin{verbatim}
if (Dq < 22612) {              // campo débil → HS es el suelo
    Eb  = -4  * Dq;            //  E(HS) = -4 Dq
    Eex = -24 * Dq + 45224.16; //  E(LS) = -24 Dq + P
    ...
} else {                       // campo fuerte → LS es el suelo
    Eb  = -24 * Dq + 45224.16;
    Eex = -4  * Dq;
}
\end{verbatim}

\textbf{Metropolis a $T$ finita}: propone voltear el espín, calcula $\Delta E$ y
acepta con
\begin{equation}
p_{\omega} = \exp\!\left(-\Delta E/(T \cdot 0.695)\right)
\end{equation}
donde $0.695$ es $\kb$ en cm$^{-1}$/K. \textbf{Barrido de $\Dq$} de $21100$ a
$24030$ (paso 1) subiendo y bajando $\Rightarrow$ lazo espín vs $\Dq$ en
\texttt{table.dat}.

\subsection{De dónde salen los números}

\begin{itemize}
\item Las rectas $-4\Dq$ y $-24\Dq+P$ son la \textbf{linealización de
  Tanabe--Sugano} para $d^6$.
\item El cruce real (de las fórmulas) es $20\Dq = 45224.16 \Rightarrow
  \Dq = 2261.2$ cm$^{-1}$, es decir $\Delta_o = 22612$ cm$^{-1}$. El umbral
  \texttt{22612} es por tanto $\Delta_o$ y \emph{no} $\Dq$: se mezclaron las dos
  magnitudes (factor 10).
\item $45224.16$ es la energía de apareamiento efectiva $P$.
\item $B=23002$ no coincide con el $B$ de Racah real del Fe(II) ($\sim$900--1050
  cm$^{-1}$); está declarado pero \textbf{no se usa}.
\end{itemize}

\subsection{Bugs y advertencias}

\begin{enumerate}
\item \textbf{$\kb$ al revés.} \texttt{exp(-dE/T * 0.695)} se evalúa como
  $(-\Delta E/T)\cdot 0.695$ (el $0.695$ multiplica en vez de dividir). Debería
  ser $-\Delta E/(T\cdot 0.695)$. Es el \emph{``error de probabilidades''} de la
  nota \texttt{notas}.
\item \textbf{\texttt{metropolis()} se llama dos veces} en \texttt{montecarlo()},
  de modo que $\Delta E$ y la aceptación no son la misma propuesta.
\item \textbf{$\Dq$ vs $\Delta_o$ (factor 10).} El umbral \texttt{22612}
  corresponde a $\Delta_o$. Las pendientes $-4/-24$ están pensadas para $\Dq$.
  El cruce queda bien localizado, pero las pendientes de energía quedan $10\times$
  desviadas.
\item \textbf{$B=23002$ muerto} (no usado).
\item \textbf{Unidades.} $\Dq\sim 22000$ cm$^{-1}$ contra $\kb T \sim 0.7$
  cm$^{-1}$ ($T=1$ K) $\Rightarrow$ el factor de Boltzmann es $0/1$ casi siempre
  (función escalón), sin ensanchamiento térmico real.
\end{enumerate}

\section{\texttt{project\_sco}: red 3D elástica / Ising}

Modelo \textbf{cooperativo} completo. Red cúbica de $N=n^3$ centros, cada uno con
posición $(x,y,z)$, espín \texttt{tSpin} ($0/1$) y \textbf{radio} dependiente del
espín: $R_{\HS}=0.418$, $R_{\LS}=0.38$, $R_{\HS}/R_{\LS}=1.1$.

Hamiltoniano (caso 1 de \texttt{f\_totalEnergy\_totalMoment}):
\begin{align}
H &= H_o + H_{nn} + H_{\mathrm{mag}}\\
H_o &= D \sum_i \sigma_i \qquad \text{(energía onsite de cambio de espín)}\\
H_{nn} &= \frac{k_{nn}}{2}\sum_{(i,j)}\bigl(r_{ij}-(R_i+R_j)\bigr)^2 \quad \text{(resorte elástico)}\\
H_{\mathrm{mag}} &= -j_{\mathrm{isin}}\sum_{(i,j)} \sigma_i\,\sigma_j \qquad \text{(acoplamiento Ising)}
\end{align}

La \textbf{cooperatividad} viene del resorte: si un centro pasa a $\HS$ crece su
$R$ y estira a los vecinos, favoreciendo su transición. El algoritmo construye la
red y los 6 vecinos, propone (voltear espín / desplazar $xyz$ / cambiar volumen
NPT) y acepta con $e^{-\Delta E/T}$, y barre $T$ en ciclo $\Rightarrow$ histéresis.

\emph{Bugs}: lector de \texttt{inputfile} vacío; $D$ sobreescrito en
\texttt{f\_switchEnergy} caso 1 ($D = -T/5 + 1$); mezcla de $D$ y $j_{\mathrm{isin}}$
en el caso 3; bug aritmético en el caso 5 (resta $H_{nnn,\mathrm{fin}}$ en vez de
$H_{nnn,\mathrm{ini}}$); RNG re-siembra \texttt{std::random\_device} en cada
llamada; unidades de $\kb$ inconsistentes; la clase \texttt{m\_lattice} es un stub.

\section{\texttt{one\_D}: cadena 1D}

Versión unidimensional del mismo modelo elástico/Ising: \texttt{Tatoms = 40},
ciclo térmico $T_{\star}=1 \to T_{\mathrm{end}}=150$ (paso 2),
$D_{\mathrm{parameter}} = -10$, promedios vivos y salida \texttt{E\_and\_M.dat} +
gnuplot (genera los \texttt{test*.pdf}). Es el predecesor directo de
\texttt{project\_sco} y sí llegó a generar resultados.

\section{\texttt{DMF.beta.*}: procesado de magnetometría}

No es simulación: procesa datos experimentales del magnetómetro (SQUID/PPMS).
Lee \texttt{inputfile.dat} (masa, masa molar, constante diamagnética de Pascal,
número de frecuencias) y calcula $M$ vs $H$ (\texttt{mvshFunction}), ZFCFC
(\texttt{zfcfcfFunction}) y susceptibilidad AC (\texttt{acFunction}). Las
versiones beta reflejan iteración sobre el mismo procesado.

\section{\texttt{project\_one}: GUI Qt}

Proyecto Qt (mainwindow + plots), probablemente una interfaz de
visualización/análisis sobre los modelos o datos. Prototipo.

\section{Errores comunes (checklist)}

\begin{enumerate}
\item \textbf{Unidades energía/temperatura.} Decidir una sola unidad (cm$^{-1}$ +
  K con $\kb=0.695$ cm$^{-1}$/K, o eV con $\kb=8.617\times10^{-5}$ eV/K) y usarla
  \emph{consistentemente} en $e^{-\Delta E/\kb T}$.
\item \textbf{$\Dq$ vs $\Delta_o$ ($10\Dq$).} No mezclarlas en el mismo bucle.
\item \textbf{Racah $A,B,C$.} $B(\mathrm{Fe}^{2+})\approx 917$--$1050$ cm$^{-1}$,
  $C\approx 4B$.
\item \textbf{RNG.} Sembrar \texttt{std::mt19937} \emph{una sola vez}.
\item \textbf{Energía propuesta = energía aceptada.} Calcular $\Delta E$ una sola
  vez por propuesta.
\item \textbf{Lector de entrada.} No dejar los \texttt{case} vacíos.
\end{enumerate}

\section{Referencias}

Material en \texttt{referencias/}. Lo esencial:
Levchenko--Khristov--Varyukhin (2014, \emph{Low Temp.\ Phys.} \textbf{40}, 571);
Lever, \emph{Inorganic Electronic Spectroscopy}; Kahn, \emph{Molecular Magnetism};
Boča, \emph{Theoretical Foundations of Molecular Magnetism}; Paulsen--Trautwein
(2004). Ver manifiesto en \texttt{referencias/README.md}.

\end{document}
